% Generated draft: requires mathematical and source review.
\documentclass[11pt]{article}
\usepackage[margin=25mm]{geometry}
\usepackage{fontspec}
\setmainfont{Latin Modern Roman}
\usepackage{amsmath,amssymb,mathtools,unicode-math}
\setmathfont{Latin Modern Math}
\usepackage{enumitem,needspace,xurl,hyperref}
\setlength{\parindent}{0pt}
\setlength{\parskip}{5pt}
\setlength{\emergencystretch}{3em}
\newcommand{\N}{\mathbb{N}}
\newcommand{\Z}{\mathbb{Z}}
\newcommand{\Q}{\mathbb{Q}}
\newcommand{\R}{\mathbb{R}}
\newcommand{\C}{\mathbb{C}}
\newcommand{\F}{\mathbb{F}}
\newcommand{\A}{\mathbb{A}}
\newcommand{\PP}{\mathbb{P}}
\newcommand{\E}{\mathbb{E}}
\newcommand{\eps}{\varepsilon}
\newcommand{\dd}{\,\mathrm{d}}
\begin{document}
\section*{UnsolvedMath ID: 701560}
\textbf{Source problem code:} catalogue entry
\subsection{Definitions and mathematical statement}
Let \(G\) be a finite group. A subset \(N\subseteq G\) is normal if \(gNg^{-1}=N\) for every \(g\in G\); equivalently, it is a union of conjugacy classes of \(G\). For a normal subset \(S\), define
\[
\kappa_G(S)=\bigl|\{D:D\text{ is a conjugacy class of }G,\ D\subseteq S\}\bigr|.
\]
For subsets \(N,C\subseteq G\), put \(NC=\{nc:n\in N,\ c\in C\}\). If \(N\) is normal and \(C\) is a conjugacy class, then \(NC\) is normal. Call \(G\) \emph{conjugacy-expansive} if
\[
\kappa_G(NC)\geq\kappa_G(N)
\quad\text{for every normal subset }N\subseteq G
\text{ and every conjugacy class }C\text{ of }G.
\]
The empty normal subset is permitted, with \(\kappa_G(\varnothing)=0\). A simple group means a nontrivial group with no nontrivial proper normal subgroup; the question does not restrict to nonabelian simple groups.

\textbf{Question.} Is every finite simple group conjugacy-expansive?
\subsection{Short English statement}
Must multiplying any union of conjugacy classes in a finite simple group by any one conjugacy class produce a union containing at least as many conjugacy classes as the original union?
\subsection{Sources}
\begin{itemize}
\item Local UnsolvedMath registry, supplied record with permanent ID \(701560\), source problem code \texttt{catalogue entry}, titled \emph{Kourovka Notebook Problem 19.12}; upstream revision \texttt{d371ee917ac9647fbb06cd6fc529deec6f3dda82}. The record labels the problem open; this is supplied metadata, not an independently verified assessment.
\item \emph{The Kourovka Notebook (Unsolved Problems in Group Theory)}, Problem 19.12, attributed in the supplied record to the 19th issue (2018), with a cited 21st-issue edition (2026), \texttt{arXiv:1401.0300v48}. Proposed by M. Bezerra, Z. Halasi, A. Mar\'oti, and S. Sidki. The supplied collection description names E. I. Khukhro and V. D. Mazurov as editors.
\end{itemize}
\end{document}
